\documentclass[../../main.tex]{subfiles}
\graphicspath{{../../image/}{../../figure/}}
\begin{document}

\section{强耗散}\label{section:ch16-强耗散}

我们在此开始研究带耗散源的平衡律组。可以看到，即使存在非齐次性，耗散也能为总变差小的初始值的 $BV$ 解保证整体存在性。然而，为使分析尽可能简单，我们通篇只考虑齐次双曲平衡律组
\begin{align}
\partial_tU(x,t)+\partial_xF(U(x,t))+P(U(x,t))=0.
\label{align:ch16-s02-e01}
\end{align}
我们假设原点 $U=0$ 包含在区域 $\mathscr{O}$ 内且 $P(0)=0$，从而 $U\equiv0$ 是平衡解。

由于分析在 $BV$ 空间中进行，我们必须对 $P$ 施加使其在 $L^1$ 中耗散的条件。为识别恰当的假设，我们作替换 $U=R(0)V$，其中 $R(U)$ 是以 $DF(U)$ 的一组线性无关右特征向量 $R_1(U),\cdots,R_n(U)$ 为列向量的 $n\times n$ 矩阵，并在 $0$ 附近线性化\eqref{align:ch16-s02-e01}。这给出系统
\begin{align}
\partial_tV_i(x,t)+\lambda_i(0)\partial_xV_i(x,t)+\sum_{j=1}^n A_{ij}V_j(x,t)=0,\qquad i=1,\cdots,n,
\label{align:ch16-s02-e02}
\end{align}
其中
\begin{align}
A=L(0)DP(0)R(0),
\label{align:ch16-s02-e03}
\end{align}
这里 $L(U)$ 是以 $DF(U)$ 的左特征向量 $L_1(U),\cdots,L_n(U)$ 为行向量的 $n\times n$ 矩阵，由 $L(U)R(U)=I$ 归一化。我们以 $\operatorname{sgn}V_i(x,t)$ 乘\eqref{align:ch16-s02-e02}，关于 $x$ 在 $(-\infty,\infty)$ 上积分，并对 $i=1,\cdots,n$ 求和，推知：当 $A$ 列对角占优，即
\begin{align}
A_{ii}-\sum_{j\neq i}|A_{ji}|\geqslant\mu>0,\qquad i=1,\cdots,n,
\label{align:ch16-s02-e04}
\end{align}
时，当 $t\to\infty$，\eqref{align:ch16-s02-e02}的解在 $L^1(-\infty,\infty)$ 中指数衰减到零。

应当指出，对角占优性质\eqref{align:ch16-s02-e04}是否成立，可能依赖于构造 $A$ 时所选用的右特征向量矩阵 $R(U)$。事实上，若取等价的特征向量矩阵 $\hat R(U)=R(U)K$，其中 $K$ 为某正定对角矩阵，则 $A$ 将被换成矩阵 $\hat A=K^{-1}AK$；而对角占优在这种相似变换下一般不保持。给定矩阵 $A$，存在正对角矩阵 $K$ 使 $K^{-1}AK$ 列对角占优，当且仅当矩阵 $\tilde A$（其元素 $\tilde A_{ii}=A_{ii}$，$i=1,\cdots,n$，$\tilde A_{ij}=-|A_{ij}|$，$i\neq j$）的所有特征值都具有正实部（见\cref{section:ch16-注记}所引文献）。
特别地，这类 $A$ 包含正三角矩阵以及行对角占优矩阵（由 Geršgorin 定理）。

对任意 $\tau>0$，从左以行向量 $\operatorname{sgn}X^\top$ 乘线性系统 $(I+\tau A)X=Y$，得到
\begin{align}
|(I+\tau A)^{-1}|\leqslant(1+\mu\tau)^{-1}.
\label{align:ch16-s02-e05}
\end{align}
正如我们将看到，正是这一性质为\eqref{align:ch16-s02-e01}的 Cauchy 问题诱导出整体解的存在性。

\begin{theorem}\label{theorem:ch16-1621}
考虑齐次、严格双曲的平衡律组\eqref{align:ch16-s02-e01}，其特征族要么真非线性、要么线性退化。设 $0$ 为平衡态，$P(0)=0$。假设对 $DF(0)$ 的某组特征向量，由\eqref{align:ch16-s02-e03}定义的矩阵 $A$ 列对角占优\eqref{align:ch16-s02-e04}。给定有界变差初始数据 $U_0$，满足\eqref{align:ch16-s01-e10}且 $\delta$ 充分小，并 $U_0(x)\to0$（当 $x\to-\infty$），则\eqref{align:ch16-s02-e01}、\eqref{align:ch16-s01-e02}存在整体可容 $BV$ 解 $U$。对每个固定 $t\in[0,\infty)$，$U(\cdot,t)$ 是 $(-\infty,\infty)$ 上的有界变差函数，且
\begin{align}
TV_{(-\infty,\infty)}U(\cdot,t)\leqslant c_1\delta e^{-\nu t},
\label{align:ch16-s02-e06}
\end{align}
其中 $c_1$ 与 $\nu$ 为正常数。
\end{theorem}

\textbf{证明梗概。}一个更强结果的详细技术性证明将在本章后续各节给出。尽管如此，我们在此提供\cref{theorem:ch16-1621}证明的一个纲要，以便用相对简单的术语向读者传达：算子分裂是如何工作的，以及矩阵 $A$ 的对角占优假设所起的作用。

我们借助\cref{section:ch16-cauchy问题}中勾勒的格式构造解 $U$。相容性的证明紧随原书定理 13.1.1 证明中的论证，此处不再重复。
只需建立约似解 $U_h$ 的总变差 $TV_{(-\infty,\infty)}U_h(\cdot,t)$ 的一个界，使其在 $h\downarrow0$ 的极限下给出所断言的估计\eqref{align:ch16-s02-e06}。

如同第十三章第 3 节那样，对 $r+s$ 为偶数者考虑菱形 $\Delta_s^r$，其顶点为 $(y_s^{r-1},t_s)$、$(y_{s-1}^r,t_{s-1})$、$(y_s^{r+1},t_s)$ 与 $(y_{s+1}^r,t_{s+1})$，见原书图 13.3.1。
目标是用从 $(x_{r-1},t_{s-1})$ 与 $(x_{r+1},t_{s-1})$ 发出的入射波扇 $\alpha$ 与 $\beta$ 的强度，来估计从 $(x_r,t_r)$ 发出的出射波扇 $\varepsilon$ 的强度。

由于我们的系统是齐次的，\eqref{align:ch16-s01-e04}给出 $V_s^r=\hat U_s^r=W_s^r$。按照算法的规定，
\begin{align}
\Omega(\alpha;U_s^{r-1})=\hat U_{s-1}^r,
\label{align:ch16-s02-e07}\\
\Omega(\beta;\hat U_{s-1}^r)=U_s^{r+1},
\label{align:ch16-s02-e08}\\
\Omega(\varepsilon;\hat U_s^{r-1})=\hat U_s^{r+1},
\label{align:ch16-s02-e09}
\end{align}
其中 $\Omega$ 是由原书(9.3.4)定义的波扇函数。

我们考虑在无源项时由 $\alpha$ 与 $\beta$ 的相互作用所产生的波扇 $\tilde\varepsilon$，即
\begin{align}
\Omega(\tilde\varepsilon;U_s^{r-1})=\Omega(\beta;\Omega(\alpha;U_s^{r-1}))=U_s^{r+1}.
\label{align:ch16-s02-e10}
\end{align}
由原书定理 9.9.1，
\begin{align}
\tilde\varepsilon=\alpha+\beta+O(1)\mathscr{D}(\Delta_s^r),
\label{align:ch16-s02-e11}
\end{align}
其中波相互作用项 $\mathscr{D}(\Delta_s^r)$ 由原书(13.3.2)定义。

我们进而把 $\varepsilon$ 与 $\tilde\varepsilon$ 联系起来。由于 $\Omega(0;U)=U$（对任意 $U\in\mathscr{O}$），\eqref{align:ch16-s02-e09}连同\eqref{align:ch16-s01-e03}与 $P(0)=0$ 给出
\begin{align}
\Omega(\varepsilon;U_s^{r-1})=U_s^{r+1}-\tau[P(U_s^{r+1})-P(U_s^{r-1})]+o(1)h|\varepsilon|,
\label{align:ch16-s02-e12}
\end{align}
其中 $o(1)$ 表示当 $\sup|U_h|$ 充分小时任意小的量。由原书(9.3.8)，
\begin{align}
\Omega(\varepsilon;U_s^{r-1})-\Omega(\tilde\varepsilon;U_s^{r-1})=\tilde R(\varepsilon-\tilde\varepsilon),
\label{align:ch16-s02-e13}
\end{align}
其中 $\tilde R$ 是某个接近 $DF(0)$ 的右特征向量矩阵 $R(0)$ 的矩阵。此外，由\eqref{align:ch16-s01-e03}，
\begin{align}
P(U_s^{r+1})-P(U_s^{r-1})=H[\hat U_s^{r+1}-\hat U_s^{r-1}]+\tau H[P(U_s^{r+1})-P(U_s^{r-1})],
\label{align:ch16-s02-e14}
\end{align}
其中 $H$ 是某个接近 $DP(0)$ 的矩阵。最后，由原书(9.3.8)与\eqref{align:ch16-s02-e09}，
\begin{align}
\hat U_s^{r+1}-\hat U_s^{r-1}=\Omega(\varepsilon;\hat U_s^{r-1})-\Omega(0;\hat U_s^{r-1})=\hat R\varepsilon,
\label{align:ch16-s02-e15}
\end{align}
其中 $\hat R$ 是某个接近 $R(0)$ 的矩阵。现在合并\eqref{align:ch16-s02-e10}、\eqref{align:ch16-s02-e12}、\eqref{align:ch16-s02-e13}、\eqref{align:ch16-s02-e14}与\eqref{align:ch16-s02-e15}，得到
\begin{align}
\tilde\varepsilon=[I+\tau\tilde A]\varepsilon+o(1)h|\varepsilon|,
\label{align:ch16-s02-e16}
\end{align}
其中
\begin{align}
\tilde A=\tilde R^{-1}[I-\tau H]^{-1}H\hat R
\label{align:ch16-s02-e17}
\end{align}
接近由\eqref{align:ch16-s02-e03}定义的矩阵 $A$。

由\eqref{align:ch16-s02-e11}、\eqref{align:ch16-s02-e16}与\eqref{align:ch16-s02-e05}，我们推出：只要 $\sup|U_h|$ 保持充分小，
\begin{align}
|\varepsilon|\leqslant(1-3\nu\tau)(|\alpha|+|\beta|)+c\mathscr{D}(\Delta_s^r),
\label{align:ch16-s02-e18}
\end{align}
其中 $\nu=\mu/4>0$。

由\eqref{align:ch16-s02-e11}、\eqref{align:ch16-s02-e16}与\eqref{align:ch16-s02-e18}，我们还推出
\begin{align}
|\varepsilon-(\alpha+\beta)|\leqslant ch(|\alpha|+|\beta|)+c\mathscr{D}(\Delta_s^r).
\label{align:ch16-s02-e19}
\end{align}

如同第十三章第 4 节那样，我们考虑网格曲线 $I$，并结合由原书(13.4.2)、(13.4.5)与(13.4.8)定义的泛函 $\mathscr{L}(I)$、$\mathscr{Q}(I)$ 与 $\mathscr{G}(I)$。
设 $J$ 是 $I$ 的直接后继（见原书图 13.4.1），我们可以重走第十三章第 4 节的分析，用\eqref{align:ch16-s02-e18}代替原书(13.4.4)得到
\begin{align}
\mathscr{L}(J)\leqslant\mathscr{L}(I)-3\nu\tau(|\alpha|+|\beta|)+c\mathscr{D}(\Delta_s^r),
\label{align:ch16-s02-e20}
\end{align}
再用\eqref{align:ch16-s02-e19}代替原书(13.4.7)得到
\begin{align}
\mathscr{Q}(J)-\mathscr{Q}(I)\leqslant ch\mathscr{L}(I)(|\alpha|+|\beta|)+[c\mathscr{L}(I)-1]\mathscr{D}(\Delta_s^r).
\label{align:ch16-s02-e21}
\end{align}
于是，当 $\kappa$ 充分大且 $\mathscr{L}(I)$ 充分小时，
\begin{align}
\mathscr{G}(J)\leqslant\mathscr{G}(I)-2\nu\tau(|\alpha|+|\beta|).
\label{align:ch16-s02-e22}
\end{align}

接着，对固定的 $s=0,1,2,\ldots$，考虑网格曲线 $J_s$，其顶点为所有 $r+s$ 为偶数的采样点 $(y_s^{r-1},t_s)$ 与 $(y_{s+1}^r,t_{s+1})$。然后假设 $\sup|U_h|$ 充分小，使得 $\mathscr{G}(J_{s-1})\leqslant2\mathscr{L}(J_{s-1})$，\eqref{align:ch16-s02-e22}给出
\begin{align}
\mathscr{G}(J_s)\leqslant(1-\nu\tau)\mathscr{G}(J_{s-1}).
\label{align:ch16-s02-e23}
\end{align}
因此，对任意 $t_s<t<t_{s+1}$，
\begin{align}
TV_{(-\infty,\infty)}U_h(\cdot,t)\leqslant(1-\nu\tau)^s\mathscr{G}(J_0),
\label{align:ch16-s02-e24}
\end{align}
其中总变差由 $\mathscr{L}(J_s)$ 度量。

由于 $U_h(x,t)\to0$（当 $x\to-\infty$），\eqref{align:ch16-s02-e24}的右端也控制 $\sup_{(-\infty,\infty)}|U_h(\cdot,t)|$。

在\eqref{align:ch16-s02-e24}的右端，$\mathscr{G}(J_0)$ 由 $c\,TV_{(-\infty,\infty)}U_0(\cdot)$ 控制。于是令 $h\downarrow0$，\eqref{align:ch16-s02-e24}给出\eqref{align:ch16-s02-e06}。证明完成。

诱导出对角占优矩阵 $A$ 的源，与第 5.5 节中遇到的耗散正定源地位相同：它们 readily 给出 Cauchy 问题解的整体存在性与指数衰减；但在物理应用中出现的系统里，它们很少遇到。在以下各节中，我们将看到如何为更广的一类源建立整体存在性，类似于第 5.5 节中的耗散半正定源。

\end{document}
